\subsection{The standard isomorphism classes}\hypertarget{tag-01PE}{}\label{fam:section:A.5}

\begin{theorem}[Classification of the standard family]\hypertarget{tag-01Q6}{}\label{fam:src:A.5.1}\label{fam:st:A.5.1}\label{wb:pshift:thm:classification}
\StatementRecord{fam:src:A.5.1}{proved}{01Q6}{\SeriesRef{app:rederive:fam-src-A.5.1}; \citet[Theorem~1.4, p.~3]{DuEtAl2026v1}}

Fix the original cusp-regular period solution, the clockwise peripheral marking, the prescribed periodic triangular cusp fan, the affine vectors $(\varepsilon,-\varepsilon')$, and the untwisted cusp collar. For every $c,c'\in\mathcal P$, the constructed compact complex threefolds satisfy
\[
 X_c\simeq X_{c'}\quad\Longleftrightarrow\quad c'-c\in2\mathbb Z.
\]
Every biholomorphism is over the identity of the original $\mathbb P^1$. The proper holomorphic family of their isomorphism classes has parameter $0<|w|<e^{-\pi M}$ and nonzero Kodaira--Spencer class at every point.
\end{theorem}



\begin{remark}[A rank-one nodal model]
\hypertarget{tag-IP09}{}\label{I:remark:rank-one-nodal-model}
The family of \cref{fam:src:A.5.1,fam:src:A.5.2} uses the original period
system with multiplicities $(3,4)$. Its arbitrary affine and cusp data are
classified in \cref{fam:section:B.2}, and its higher multiplicities in
\cref{fam:section:B.3}.
For the elliptic surface $y^2=x^3+tx+1$ of
\cite[Section~2, pp.~3--5; Proposition~3.1, pp.~6--7;
Proposition~6.1, pp.~17--18]{ViaclovskyTwo2026}, fix the nodal value
$p_1$ with $4p_1^3+27=0$. In that paper's equation~(3.1), put
$\gamma_0=\gamma(0)$, $\beta_0=\beta(0)$ and
$c_0=\lambda\delta(0)$, with $0<|c_0|<1$.
The normalization of this rank-one filling is
\[
 \begin{gathered}
 E'=\mathbb C^*/\langle c_0\rangle,\qquad
 R=\mathbb P_{E'}(\mathcal O_{E'}\oplus L),\\
 R=(\mathbb P^1\times\mathbb C^*)/
       \langle(z,v)\mapsto(\gamma_0z,c_0v)\rangle,
 \end{gathered}
\]
and its boundary sections are identified by $v\mapsto\beta_0v$.
The invariant nodal volume form in equation~(6.2) of that paper
trivializes the ambient canonical line near this fibre. The branch-normal
map therefore has quotient $\nu_*K_R$ by the local coefficient map in
\eqref{fam:eq:ab-normal-quotient}.
For the ruling $\pi:R\to E'$, one has
$\pi_*\mathcal O_R=\mathcal O_{E'}$ and
$R^1\pi_*\mathcal O_R=0$; Leray and Serre duality give
$h^\bullet(K_R)=(0,1,1)$.
The corresponding summand of the connecting homomorphism is
\[
 H^1(R,K_R)=\mathbb C\longrightarrow H^2(T_X(-\log D)),
\]
where $D$ is the union of that model's three reduced nodal fibres.
Its defining boundary-translation \v Cech cocycle uses the full
holomorphic units in equation~(3.1), their central values
$(\gamma_0,c_0,\beta_0)$ and the vertical frames. The integral
comparison equations use that model's four monodromy matrices and
their collar defects. These are the explicit local and global inputs
for this model; the displayed cohomology calculation evaluates the
source of the connecting map.
\end{remark}

